\documentclass[11pt]{article}

\usepackage[margin=1in]{geometry}
\usepackage[T1]{fontenc}
\usepackage[utf8]{inputenc}
\usepackage{lmodern}
\usepackage{microtype}
\usepackage{hyperref}
\usepackage{amsmath}
\usepackage{enumitem}

\title{Future Work: Sequential Localization with Prediction-Based Filtering}
\author{VPR Tutorial Project Notes}
\date{\today}

\begin{document}
\maketitle

\section*{Motivation}
The current Visual Place Recognition (VPR) workflow is mainly frame-by-frame: for each incoming image, the system finds the best matching place in the map. This works as a baseline, but it can be unstable under viewpoint and perspective changes, motion blur, and repetitive-looking scenes.

A strong future-work direction is to move from \emph{single-frame matching} to \emph{sequential state estimation}. The key idea is to combine:
\begin{itemize}[leftmargin=*]
  \item what the robot \textbf{expects} to see next (from motion and map), and
  \item what the camera \textbf{actually} observes (from VPR inference).
\end{itemize}

This is exactly the setting where a Kalman-style filter is useful.

\section*{Core Idea from the Discussion}
The professor's point can be summarized as:
\begin{quote}
``Do not localize each frame independently. Use robot motion and map geometry to predict the next pose and expected observation, then fuse that prediction with visual matching evidence over time.''
\end{quote}

So instead of only asking:
\begin{itemize}[leftmargin=*]
  \item ``Which map image best matches this frame?''
\end{itemize}
we ask:
\begin{itemize}[leftmargin=*]
  \item ``Given where I was and how I moved, where should I be now?''
  \item ``Does the camera observation agree with that prediction?''
\end{itemize}

\section*{What ``Expectation'' Means in Practice}
At time $t$, assume the robot has pose state:
\[
\mathbf{x}_t = [x_t,\; y_t,\; \theta_t,\; v_t,\; \omega_t]^\top
\]
where $(x,y)$ is 2D location, $\theta$ is heading, and $v,\omega$ are linear/angular velocity terms.

Using odometry or control input, a motion model predicts the next state:
\[
\hat{\mathbf{x}}_{t+1}^{-} = f(\mathbf{x}_t,\mathbf{u}_t)
\]
This is the \textbf{expectation}: where the robot likely is before seeing the next camera frame.

From this predicted pose and the 2D map, the system can also infer what place/region should be visible.

\section*{What ``Observation'' Means}
At time $t+1$, the camera frame is processed by VPR to produce evidence such as:
\begin{itemize}[leftmargin=*]
  \item best matching map index,
  \item similarity score,
  \item top-$k$ candidate matches.
\end{itemize}

This acts as a noisy measurement $\mathbf{z}_{t+1}$ of location.

\section*{Why Kalman-Style Filtering Helps}
A Kalman filter framework alternates two steps:
\begin{enumerate}[leftmargin=*]
  \item \textbf{Prediction:} propagate state using robot motion model.
  \item \textbf{Update:} correct predicted state using visual measurement.
\end{enumerate}

It also propagates uncertainty, so the system can decide how much to trust motion vs vision at each step.

Benefits for this project:
\begin{itemize}[leftmargin=*]
  \item \textbf{Temporal consistency:} fewer sudden pose jumps between unrelated map places.
  \item \textbf{Robustness to bad frames:} one poor visual match has limited impact.
  \item \textbf{Better navigation behavior:} smoother and physically plausible trajectories.
  \item \textbf{Perspective-change resilience:} sequence context compensates for viewpoint variation.
\end{itemize}

\section*{Conceptual Model for Future Implementation}
A practical formulation would use:
\begin{itemize}[leftmargin=*]
  \item \textbf{Process model} (motion): predicts $(x,y,\theta)$ over time,
  \item \textbf{Measurement model} (vision): maps VPR match evidence to pose observation,
  \item \textbf{Fusion layer} (filter): combines both using uncertainty.
\end{itemize}

At each time step:
\begin{align*}
\hat{\mathbf{x}}_{t}^{-} &= f(\hat{\mathbf{x}}_{t-1}, \mathbf{u}_{t-1}) \\
\hat{\mathbf{x}}_{t} &= \hat{\mathbf{x}}_{t}^{-} + \mathbf{K}_t\big(\mathbf{z}_t - h(\hat{\mathbf{x}}_{t}^{-})\big)
\end{align*}
where $\mathbf{K}_t$ is the Kalman gain and $h(\cdot)$ maps state to expected measurement.

\section*{Important Note on Filter Choice}
A standard linear Kalman filter assumes linear dynamics and Gaussian noise. In real VPR localization:
\begin{itemize}[leftmargin=*]
  \item pose and camera observation relationships are often nonlinear,
  \item match hypotheses may be multi-modal in ambiguous scenes.
\end{itemize}

So likely candidates are:
\begin{itemize}[leftmargin=*]
  \item \textbf{EKF} (Extended Kalman Filter) as a strong baseline,
  \item \textbf{UKF} (Unscented Kalman Filter) for stronger nonlinear behavior,
  \item \textbf{Particle filter} if ambiguity is highly multi-modal.
\end{itemize}

\section*{How to Position This as Future Work}
A concise statement for report/presentation:
\begin{quote}
``Future work will integrate VPR with probabilistic sequential localization in a 2D mapped environment. Motion/odometry will provide state prediction, while VPR similarity will provide measurement updates. We will evaluate gains in trajectory stability, relocalization accuracy, and robustness under viewpoint and perspective changes.''
\end{quote}

\section*{Expected Evaluation Criteria}
When this is implemented, useful evaluation metrics include:
\begin{itemize}[leftmargin=*]
  \item localization error over trajectory (position and heading),
  \item temporal smoothness / trajectory consistency,
  \item relocalization robustness after visual degradation,
  \item precision-recall style place-recognition metrics before and after filtering.
\end{itemize}

\end{document}
